\PBchapter{Infância e adolescência}

%\PByear{1977}

Como já mencionei, nasci em Abril de 1977 na cidade de Cascavel, oeste do Paraná.

Aos sete anos de idade, meus pais se separaram. Naquela época, eu morava em Foz do Iguaçu, no oeste do Paraná — uma cidade de fronteira, onde Brasil, Paraguai e Argentina praticamente se encontram.

E, olhando hoje, não consigo deixar de pensar: será que a história dos meus avós e dos meus pais, antes de mim, já não indicava morar nesta região?

A separação aconteceu. 

E, como na maioria das vezes, não veio acompanhada de um plano claro.

Vieram mudanças.

Ajustes.

E, principalmente, a necessidade de adaptação.

No fim dos anos 80, o Brasil também não oferecia muita estabilidade. A inflação fazia parte do cotidiano. Planos econômicos surgiam e desapareciam. Tudo parecia provisório.

De certa forma, a gente crescia no meio disso — aprendendo a lidar com o incerto sem nem perceber.

Em 1987, retornamos para Cascavel. Minha mãe precisava fazer uma cirurgia, e a ideia era que minha vó cuidasse de mim e da minha irmã.

%\PByear{1989}

Por volta dos 12 anos, em 1989, comecei a entender que a vida não era tão simples quanto parecia.

Muitas perguntas surgiam. 

Muitos medos. 

Ouvia muitas conversas, baixas, mas dava para ouvir. 

Coisas como... 

E se a mãe deles morrer? 

O que vai ser desses meninos? 

Vamos ter que mandar de volta para o pai?

Aos 14, essa sensação só aumentou.

E, muitas vezes, as respostas simplesmente não vinham.

Já na adolescência, decidi fazer um curso de Mecânica Geral no Senai. 

A ideia, na época, era simples. 

Seguir algo próximo do que meu pai fazia. 

Talvez repetir um caminho que, para mim, parecia lógico.

Mas não foi bem assim.

Durante o curso, percebi que a formação era voltada para torneiro mecânico. 

Ainda assim, segui até o fim. 

Fiz estágio durante o período.

E foi só.

Nunca atuei na área.

Foi uma das primeiras vezes em que percebi, na prática, que aquilo que parece fazer sentido no início nem sempre se sustenta ao longo do caminho.

%\PByear{1994}

Algum tempo depois, em 1994, já no ensino médio, minha mãe ficou desempregada. 

Surgiu então a possibilidade e ela foi para Umuarama, em busca de trabalho.

Enquanto isso, permaneci em Cascavel, cursando o segundo ano e tentando seguir com o que era possível naquele momento.

Foi nessa fase que consegui meu primeiro trabalho, real, pois antes eram coisas como, desentortar prego, servente de pedreiro, ajudar com mudança, babá, limpeza de escritório.

O trabalho era em uma distribuidora de revistas.

Nada planejado. 

Surgiu por indicação de um colega de escola.

Minha mãe acabou retornando de Umuarama.

Eu já estava trabalhando. 

Nem tinha contado para ela.

E, pouco tempo depois, ela também conseguiu um emprego — novamente por meio de contatos. 

Algo que, visto hoje, parece simples, mas que na época fazia toda a diferença.

Era o período do Plano Real.

Uma tentativa de estabilidade em meio a tantos anos de incerteza.

O Brasil comemorava o título da Copa do Mundo de 1994.

Chorava a perda de um ídolo, Airton Senna. 

Enquanto, na vida real, cada família buscava seu próprio equilíbrio.

Fiquei neste trabalho por um bom tempo.

Não era um sonho.

Mas era o que havia.

E, muitas vezes, “o que há” é exatamente o que nos sustenta.

Com o tempo, aquele primeiro trabalho deixou de ser apenas uma novidade.

Virou rotina.

E, como acontece com muitas coisas na vida, a rotina vai tirando o peso das escolhas e transformando tudo em algo automático.

Eu trabalhava, estudava e seguia.

Sem grandes planos.

Sem uma visão clara de futuro.

Era mais sobre continuar do que sobre escolher.

Em algum momento, após concluir o ensino médio, estudar deixou de ser prioridade.

Não por decisão consciente, mas por consequência. 

O trabalho ocupava espaço. 

A necessidade falava mais alto.

E, quando se percebe, já se está vivendo no modo “seguir em frente”, sem muito tempo para questionar.

Foi uma fase de aprendizados — não necessariamente aqueles que se encontram em livros ou salas de aula.

Aprendi sobre responsabilidade.

Sobre compromisso.

E, principalmente, sobre o peso de ter que fazer as coisas acontecerem, mesmo sem entender exatamente para onde se está indo.

As oportunidades que surgiam nem sempre tinham relação com algum plano maior.

E talvez esse seja o ponto.

Não havia plano.

Havia necessidade.

E isso molda decisões.

Molda caminhos.

Molda, muitas vezes, a própria forma de enxergar a vida.

Com o passar do tempo, novas experiências profissionais foram surgindo. Algumas melhores, outras nem tanto. 

Mas todas contribuíram de alguma forma.

Nem sempre era possível escolher.

Na maioria das vezes, era preciso aceitar.

E seguir.

O curioso é que, olhando para trás, tudo parece fazer algum sentido — mesmo aquilo que, na época, parecia completamente desconectado.

Mas essa percepção só vem depois.

Naquele momento, era apenas a vida acontecendo.

Sem roteiro.

Sem garantias.

E, muitas vezes, sem espaço para parar e pensar.

Talvez seja por isso que, hoje, quando tento encontrar um “propósito”, eu não encontre.

Porque, durante muito tempo, não houve espaço para isso.

Houve movimento.

Houve tentativa.

Houve continuidade.

Mas não houve pausa.

E sem pausa, algumas perguntas simplesmente não aparecem.

Ou aparecem tarde demais.

Com o passar dos anos, a vida foi ganhando uma espécie de forma.

Não uma forma planejada.

Mas a forma possível.

Entre trabalhos, mudanças e tentativas de acerto, fui construindo um caminho. 

Não aquele que a gente imagina quando é mais novo, cheio de certezas e objetivos claros.

Ser bombeiro.

Mas um caminho real.

Feito de escolhas possíveis.

De decisões tomadas mais pela necessidade do que pela convicção.

E, ainda assim, um caminho.

Em alguns momentos, as coisas davam certo.

Em outros, nem tanto.

Mas a vida não parava para avaliação.

Ela simplesmente seguia.

E eu seguia junto.

O tempo foi passando quase sem pedir licença.

E, quando se percebe, já não se está mais no começo.

Já existe uma história.

Já existem marcas.

Algumas conquistas.

Alguns arrependimentos.

E uma série de episódios que, isoladamente, podem parecer pequenos, mas que, juntos, formam quem a gente se torna.

Talvez o ponto mais honesto seja esse: não existe uma grande virada.

Não existe um momento único, épico, que define tudo.

Existe uma sequência.

Dias comuns.

Decisões simples.

Caminhos que vão se formando sem anúncio.

E, no meio disso tudo, uma sensação que às vezes incomoda.

A de que poderia ter sido diferente.

Mas, ao mesmo tempo, a consciência de que foi exatamente assim.

Sem glamour.

Sem roteiro perfeito.

Sem propósito claramente definido.

E talvez seja isso que mais me chama atenção hoje.

Não o que faltou.

Mas o fato de que, mesmo sem todas as respostas, a vida aconteceu.

E continua acontecendo.

Talvez o propósito nunca tenha sido algo a ser encontrado.

Talvez tenha sido apenas seguir.

Mesmo sem entender completamente o porquê.

Essas reflexões todas vieram, e eu não tinha nem 19 anos.

Não tinha com quem conversar sobre isso.

Não tinha uma referência.

E assim foi minha fase de infância até a adolescência.


%\PBreflection{É perfeitamente plaúsivel, que a vida não vem com manual de instrução}